\documentclass[11pt]{article}
\usepackage[T1]{fontenc}
\usepackage[utf8]{inputenc}
\usepackage{lmodern}
\usepackage[margin=1in]{geometry}
\usepackage{amsmath,amssymb,amsthm}
\usepackage{booktabs,longtable}
\usepackage[title]{appendix}
\usepackage[colorlinks=true,linkcolor=blue,citecolor=blue,urlcolor=blue]{hyperref}
\hypersetup{pdfauthor={P. Graham},pdftitle={A lattice parallelepiped with no unimodular triangulation}}
\newtheorem{theorem}{Theorem}[section]
\newtheorem{lemma}[theorem]{Lemma}
\newtheorem{proposition}[theorem]{Proposition}
\newtheorem{corollary}[theorem]{Corollary}
\theoremstyle{definition}
\newtheorem{question}[theorem]{Question}
\newtheorem{remark}[theorem]{Remark}
\newcommand{\Z}{\mathbb{Z}}
\newcommand{\R}{\mathbb{R}}
\newcommand{\conv}{\operatorname{conv}}
\newcommand{\relint}{\operatorname{relint}}
\title{A lattice parallelepiped with no unimodular triangulation}
\author{P. Graham\\ \small\url{https://orcid.org/0009-0007-8493-1505}}
\date{November 5, 2026}
\begin{document}
\maketitle
\begin{abstract}
A lattice polytope has a unimodular triangulation if it can be cut into
lattice simplices of the smallest possible volume. Haase, Paffenholz,
Piechnik and Santos asked in their 2017 survey whether every lattice
parallelepiped, and more generally every lattice zonotope, has one. Such
polytopes have the integer decomposition property, and Codenotti and Santos
proved that every 3-dimensional lattice parallelepiped can be covered by
unimodular simplices. We show that the answer is no in every dimension $d\ge 4$.
The 4-dimensional parallelepiped $P$ spanned by $e_1,e_2,e_3$ and
$(1,2,3,5)$ has no unimodular triangulation, and its products with unit
cubes give examples in higher dimensions. Its edges at the origin span
a cone that Firla and Ziegler showed cannot be triangulated into unimodular
cones generated by its Hilbert basis (its minimal generators). Yet the
cone of $P$ at each vertex can be so triangulated with lattice points of
$P$: the obstruction is not seen at any vertex. The
proof is by hand: points of $P$ with coordinates in $\frac12\Z$ and
$\frac14\Z$ force faces into any unimodular triangulation, until one such
point lies in none. Checked certificates show that $P$ has a unimodular
cover and is, up to lattice equivalence, the unique 4-dimensional example of
determinant at most 5. The question in
dimension 3 remains open.

This work has not yet been reviewed by experts. The proofs, the checking programs and the text were produced by AI models (from OpenAI and Anthropic), working under the author's direction.
\end{abstract}
\medskip
\noindent\textbf{MSC 2020:} 52B20, 52B11, 14M25.\\
\textbf{Keywords:} unimodular triangulation, lattice parallelepiped, lattice zonotope, unimodular cover, IDP polytope.

\section{Introduction}

A \emph{lattice polytope} $P\subset\R^d$ is the convex hull of finitely many
points of $\Z^d$. A lattice $d$-simplex is \emph{unimodular} if its vertices
$v_0,\dots,v_d$ satisfy $|\det(v_1-v_0,\dots,v_d-v_0)|=1$, that is, if the
edge vectors at one vertex form a basis of $\Z^d$; equivalently, its volume
is $1/d!$. A \emph{unimodular triangulation} of $P$ is a triangulation of $P$
into unimodular simplices: a finite set of them whose union is $P$, any two
of which meet in a common face (possibly empty). A \emph{unimodular cover}
is a finite set of unimodular simplices whose union is $P$, with no
condition on how they meet. $P$ has the \emph{integer decomposition
property} (is \emph{IDP}, also called \emph{normal}) if every lattice point
of $kP$ is a sum of $k$ lattice points of $P$, for every positive integer
$k$. Each property implies the next:
\[
\text{unimodular triangulation}\ \Rightarrow\ \text{unimodular cover}\ \Rightarrow\ \text{IDP}.
\]
In their survey, Haase, Paffenholz, Piechnik and Santos list these and
intermediate properties in a numbered chain that includes the three above,
and state that ``All the implications (i) $\Rightarrow$ (i+1) are
strict''~\cite[p.~10]{HPPS}, citing Bruns and Gubeladze~\cite{BG}. In a
section first added in the third arXiv version (December 11, 2017), they
ask~\cite[\S1.5.4]{HPPS}:
\begin{quote}
Lattice parallelepipeds and, more generally, lattice zonotopes are known to
be IDP. Is it true that they all have unimodular triangulations?
\end{quote}
A \emph{zonotope} is a Minkowski sum of segments; a parallelepiped is one
with $d$ linearly independent segments in $\R^d$. Ferroni and Higashitani
repeat that ``it is not known whether zonotopes admit unimodular
triangulations in general''~\cite[\S6.1]{FH}.

\paragraph{Notation.}
For an integer matrix $B$ with $\det B\neq0$ we write
$P(B)=\{Bt: t\in[0,1]^d\}$ for the lattice parallelepiped spanned by the
columns of $B$, and, in dimension $4$, $P[a,b,c,n]$ for $P(B)$ with columns
$e_1,e_2,e_3$ and $(a,b,c,n)$. Its normalized volume is $d!\,|\det B|$; we
call $|\det B|$ its \emph{determinant}. Two parallelepipeds are
\emph{equivalent} if an affine map $x\mapsto Ux+u$ with $U\in GL_d(\Z)$ and
$u\in\Z^d$ carries one onto the other; equivalent lattice polytopes have
unimodular triangulations or covers together.

\subsection{Main results}

\begin{theorem}\label{thm:main}
The lattice parallelepiped $P=P[1,2,3,5]$, spanned by $e_1,e_2,e_3$ and
$(1,2,3,5)$, has no unimodular triangulation.
\end{theorem}

\begin{corollary}\label{cor:all}
For every $d\ge4$ the $d$-dimensional lattice parallelepiped
$P\times[0,1]^{d-4}$ has no unimodular triangulation. Hence, in every
dimension $d\ge4$, neither all lattice parallelepipeds nor all lattice
zonotopes have unimodular triangulations.
\end{corollary}

Theorems~\ref{thm:cover}
and~\ref{thm:classes} and Propositions~\ref{prop:cones}
and~\ref{prop:family} rest on exact computations with checked certificates
(Section~\ref{sec:verification}); Theorem~\ref{thm:main} and
Corollary~\ref{cor:all} do not.

\begin{theorem}\label{thm:cover}
$P$ has a unimodular cover.
\end{theorem}

\begin{theorem}\label{thm:classes}
Up to equivalence, the $4$-dimensional lattice parallelepipeds of
determinant at most $8$ without a unimodular triangulation are exactly four
classes: that of $P$ (which contains $P[1,2,2,5]$), that of $P[1,2,2,7]$
(which contains $P[1,4,3,7]$; see Section~\ref{sec:prior}), and those of
$P[1,2,3,8]$ and $P[1,3,3,8]$. In particular, $P$ is, up to equivalence,
the unique $4$-dimensional lattice parallelepiped of determinant at most $5$
without a unimodular triangulation.
\end{theorem}

A pointed rational cone $C\subset\R^4$ has a \emph{Hilbert basis}, the
minimal set of lattice points generating the semigroup $C\cap\Z^4$; a cone
is \emph{unimodular} if it is generated by a basis of $\Z^4$. The
\emph{vertex cone} of $P$ at a vertex $v$ is the cone generated by $P-v$.

\begin{proposition}\label{prop:cones}
At each of the $16$ vertices $v$ of $P$, the vertex cone has a
triangulation into unimodular cones whose rays pass through lattice points
of $P-v$. At $v=0$ one such triangulation has $56$ cones and uses as rays all
$19$ nonzero lattice points of $P$, of which $11$ are not in the Hilbert
basis of the cone.
\end{proposition}

\begin{proposition}\label{prop:family}
For every odd $p$ with $5\le p\le 31$ the parallelepiped $P[1,2,2,p]$ has no
unimodular triangulation.
\end{proposition}

\subsection{Prior work}\label{sec:prior}

Codenotti and Santos proved that ``Every 3-dimensional lattice
parallelepiped has a unimodular cover''~\cite[Theorem~1.2]{CS} and asked
\cite[Questions~1.4 and~1.5]{CS}: ``Do 3-dimensional parallelepipeds have
unimodular triangulations?'' and ``Higher dimensional parallelotopes
(affine images of cubes) are IDP. Do they have unimodular covers?'' After
stating their cover results (Theorems~1.2 and~1.6 and Corollaries~1.3, 1.7
and~1.8), they add: ``We believe that the 3-polytopes in all these
statements have unimodular triangulations, but this remains an open
question.'' Corollary~\ref{cor:all} does not concern dimension~3, and by
Theorem~\ref{thm:cover} $P$ does not answer Question~1.5 of~\cite{CS}.

Our example comes from the theory of Hilbert bases of cones. A
\emph{Hilbert triangulation} of a cone is a triangulation into unimodular
cones generated by elements of its Hilbert basis; Firla and Ziegler call it
a Hilbert Partition~\cite[Proposition~3]{FZ}. Bruns, Gubeladze and Trung
record that Sebő's theorem~\cite{Se} that every 3-dimensional cone has a
unimodular Hilbert triangulation ``cannot be extended to dimension $\ge 4$,
as shown by a counterexample due to Bouvier and
Gonzalez-Sprinberg''~\cite[Section~3.1 and Theorem~3.1.3]{BGT}. That
counterexample is the cone spanned by $e_1,e_2,e_3$ and
$(1,3,4,7)$~\cite[Exemple~3.1]{BGS}; its parallelepiped is equivalent to
$P[1,4,3,7]$, the exception of determinant 7 in Theorem~\ref{thm:classes}.
Writing $C[a_1,\dots,a_n]$ for the cone spanned by the first $n-1$ unit
vectors and $(a_1,\dots,a_n)$, Firla and Ziegler write: ``The saga of
counterexamples in this field starts with $C[1,4,3,7]$, the example of
Bouvier \& Gonzalez-Sprinberg \dots{} from 1992, of a 4-dimensional cone
that does not have a Hilbert Partition. A slight simplification yields the
following.'' Their Example~8 follows: ``The 4-dimensional cone
$C[1,2,3,5]$ does not have a Hilbert Partition''~\cite{FZ}. This is the
vertex cone of $P$ at the origin.

Bruns, Gubeladze and Trung also state that ``There exists a normal
4-dimensional lattice polytope \dots{} that has no unimodular
triangulation''~\cite[Proposition~1.2.4(c)]{BGT97}. Their polytope is the
convex hull of $0$ and the Hilbert basis $E$ of the
Bouvier--Gonzalez-Sprinberg cone; its only lattice points are $0$ and the
points of $E$, and at the vertex $0$ a unimodular triangulation of it would
give a unimodular Hilbert triangulation of the cone~\cite[proof of
Proposition~1.2.4(c), after Theorem~1.3.1]{BGT97}. That argument does not
apply to $P$: a triangulation of $P$ may use lattice points of $P$ outside
the Hilbert basis, and by Proposition~\ref{prop:cones} every vertex cone of
$P$ can be triangulated unimodularly with lattice points of $P$.
Theorem~\ref{thm:main} gives a normal polytope without a unimodular
triangulation that is a parallelepiped, hence a zonotope.
Appendix~\ref{app:unread} lists the sources we could not read in
full.

\section{Proof overview}

The obstruction in $P$ is not at a vertex (Proposition~\ref{prop:cones}),
and the step that needs an idea is to find where it sits. Suppose $T$ is a
unimodular triangulation of $P$. Every point $y\in P$ lies in the relative
interior of exactly one face $G$ of $T$, and since $G$ is unimodular, a
point with $Ny\in\Z^4$ has barycentric coordinates in $G$ that are positive
multiples of $1/N$ (Lemma~\ref{lem:den}). For $N=2$ this forces a face to
be an edge with $y$ as its midpoint, and for most half-integral points of
$P$ there is only one pair of lattice points of $P$ with that midpoint: we
get twelve edges of $T$ for free. For $N=4$ there are finitely many
candidate faces, listed by hand; two faces of $T$ meet in a common face
(Lemma~\ref{lem:common}), and this rules out every candidate that crosses a
known face. One quarter-integral point $q$ then forces a tetrahedron on four
lattice points near the origin, one of whose edges passes through the
center of $P$, to be a face of $T$, and a second quarter-integral point $r$
then lies in the relative interior of no possible face. We work in the
equivalent parallelepiped $P[1,2,2,5]$, whose lattice points have last
coordinate $0$ to $5$; this coordinate organizes every case split.

\section{Two lemmas}\label{sec:lemmas}

In this section $Q$ is a lattice polytope of dimension $d$ and $T$ is a
unimodular triangulation of $Q$. A \emph{face of $T$} is a face of one of its
simplices, vertices included.

\begin{lemma}\label{lem:den}
Let $N\ge1$ and let $y\in Q$ with $Ny\in\Z^d$. Then $y$ lies in the relative
interior of exactly one face $G$ of $T$, and $y=(m_1g_1+\dots+m_kg_k)/N$,
where $g_1,\dots,g_k$ are the vertices of $G$ and $m_1,\dots,m_k$ are
positive integers with sum $N$. In particular, every lattice point of $Q$
is a vertex of $T$, and a point $y$ with $2y\in\Z^d\setminus 2\Z^d$ is the
midpoint of an edge of $T$.
\end{lemma}

\begin{proof}
Since $T$ covers $Q$, $y$ lies in a simplex $\sigma=\conv(v_0,\dots,v_d)$ of
$T$; let $G$ be the smallest face of $\sigma$ containing $y$, so that
$y\in\relint G$. Because $\sigma$ is unimodular, $v_1-v_0,\dots,v_d-v_0$ is a
basis of $\Z^d$, so every lattice point has integer barycentric coordinates
with respect to $\sigma$. If $\lambda$ are the barycentric coordinates of $y$,
then the lattice point $Ny-(N-1)v_0$ has coordinates
$N\lambda-(N-1)(1,0,\dots,0)$, so every $N\lambda_i$ is an integer, and
$\lambda_i>0$ exactly on the vertices of $G$; put $m_i=N\lambda_i$. If also
$y\in\relint G'$ for a face $G'$ of $\sigma'\in T$, then $y$ lies in the
common face $F=\sigma\cap\sigma'$, which contains $G$ and $G'$ as the smallest
faces of $\sigma$ and $\sigma'$ containing $y$; in the simplex $F$ only one
face contains $y$ in its relative interior, so $G=G'$.
\end{proof}

\begin{lemma}\label{lem:common}
Let $\conv(S)$ and $\conv(R)$ be faces of $T$ with vertex sets $S$ and $R$,
and suppose that
\[
\sum_{s\in S}p_s\,s=\sum_{r\in R}q_r\,r\quad\text{with integers } p_s,q_r\ge0
\text{ and } \sum_s p_s=\sum_r q_r>0.
\]
Then $p_s=0$ for every $s\in S\setminus R$.
\end{lemma}

\begin{proof}
Let $\conv(S)\subseteq\sigma_1$ and $\conv(R)\subseteq\sigma_2$ with
$\sigma_1,\sigma_2\in T$, and let $F=\sigma_1\cap\sigma_2$, a common face.
Then $\conv(S)\cap F$ and $\conv(R)\cap F$ are faces of the simplex $F$,
spanned by the vertices of $S$ and of $R$ that lie in $F$, and their
intersection is the face of $F$ spanned by their common vertices. Hence
$\conv(S)\cap\conv(R)=\conv(S\cap R)$. The point
$z=\sum_s p_s s/\sum_s p_s$ lies in both, hence in $\conv(S\cap R)$. Since
$S$ is affinely independent, the barycentric coordinates $p_s/\sum p_s$ of
$z$ in $\conv(S)$ are unique, so they vanish outside $S\cap R$.
\end{proof}

We use Lemma~\ref{lem:common} to rule out a candidate face $\conv(S)$: we
exhibit an identity of this form, with $R$ already known to be a face of $T$,
in which some $s\notin R$ has $p_s>0$. When $R$ is a single lattice point
$x\notin S$, this says that $\conv(S)$ contains $x$, while a unimodular
simplex contains no lattice points other than its vertices.

\section{Proof of Theorem~\ref{thm:main}}\label{sec:proof}

\paragraph{An equivalent parallelepiped.}
Let $U\in GL_4(\Z)$ have rows $(1,0,0,0)$, $(0,1,0,0)$, $(0,0,-1,1)$,
$(0,0,0,1)$. It maps $e_1,e_2,e_3,(1,2,3,5)$ to $e_1,e_2,-e_3,(1,2,2,5)$, so
$x\mapsto Ux+e_3$ maps $P[1,2,3,5]$ onto $P'=P[1,2,2,5]$, and it suffices to
show that $P'$ has no unimodular triangulation. Writing
$x=t_1e_1+t_2e_2+t_3e_3+t_4(1,2,2,5)$ gives $t_4=x_4/5$, $t_1=x_1-x_4/5$ and
$t_i=x_i-2x_4/5$ for $i=2,3$, so $x\in P'$ if and only if
\[
0\le x_4\le5,\qquad \tfrac{x_4}{5}\le x_1\le1+\tfrac{x_4}{5},\qquad
\tfrac{2x_4}{5}\le x_2,x_3\le 1+\tfrac{2x_4}{5}.
\]
For integer $x_4=k$ this allows $x_1\in\{0,1\}$ and $x_2,x_3\in\{0,1\}$ if
$k=0$; $x=(1,1,1,k)$ if $k=1,2$; $x=(1,2,2,k)$ if $k=3,4$; and
$x_1\in\{1,2\}$, $x_2,x_3\in\{2,3\}$ if $k=5$. So $P'$ has exactly $20$
lattice points. With $a,b\in\{0,1\}$ we name them
\[
A_{ab}=(0,a,b,0),\quad B_{ab}=(1,a,b,0),\quad C_{ab}=(1,2+a,2+b,5),\quad D_{ab}=(2,2+a,2+b,5),
\]
the $16$ vertices of $P'$, and the interior points $I_1=(1,1,1,1)$,
$I_2=(1,1,1,2)$, $I_3=(1,2,2,3)$, $I_4=(1,2,2,4)$; the last coordinate of
$I_k$ is $k$. For example $A_{01}=(0,0,1,0)$ and $C_{10}=(1,3,2,5)$.
Suppose, for a contradiction, that $T$ is a unimodular triangulation of
$P'$.

\paragraph{Step 1: twelve forced edges.}
Each row of Table~\ref{tab:step1} gives lattice points $u,w$ of $P'$ whose
sum $2y=u+w$ has an odd coordinate. Then $y\in P'$ is not a lattice point,
so by Lemma~\ref{lem:den} it is the midpoint of an edge of $T$, whose
endpoints are two distinct lattice points of $P'$ with sum $2y$. Their last
coordinates add up to the last coordinate of $2y$, which is at most $7$,
and the table runs through every way to split it; only $u+w$ gives the sum
$2y$. Hence each of the twelve segments $uw$ in Table~\ref{tab:step1} is an
edge of $T$.

\paragraph{Step 2: the tetrahedron $A_{11}I_1I_2I_3$ is a face of $T$.}
Let $q=(A_{11}+I_1+I_2+I_3)/4$, so that $4q=(3,5,5,6)$; then $q\in P'$ and $q$
is not a lattice point. By Lemma~\ref{lem:den} the face of $T$ that contains
$q$ in its relative interior is $\conv(S)$, where $S$ is a set of lattice
points of $P'$ carrying positive integer weights with sum $4$ and weighted
sum $(3,5,5,6)$. We write a weighted set as a sum: $A_{00}+B_{11}+2I_3$ is
$\{A_{00},B_{11},I_3\}$ with weights $1,1,2$; and $A_{11}\,I_2$ is the face
with vertices $A_{11}$ and $I_2$. There are exactly $21$ such weighted sets
(Lemma~\ref{lem:enum}). For each of the $20$ other than
$A_{11}+I_1+I_2+I_3$, Table~\ref{tab:step2} gives an identity of
Lemma~\ref{lem:common} with a known face $R$, either an edge from Step~1 or
a lattice point. The coefficients of an identity need not be the weights of
$S$: in the first row $S=A_{00}+B_{11}+2I_3$, and the identity
$2A_{00}+B_{11}+2I_3 = 2A_{11}+3I_2$ (both sides have total weight $5$) has
$p_{A_{00}}=2>0$ with $A_{00}\notin\{A_{11},I_2\}$, so
$\conv(A_{00},B_{11},I_3)$ is not a face of $T$. Hence
$\conv(A_{11},I_1,I_2,I_3)$ is a face of $T$, and so is its edge $I_2I_3$,
whose midpoint is the center $(1,\tfrac32,\tfrac32,\tfrac52)$ of $P'$.

\paragraph{Step 3: a point that no face can contain.}
Let $r=(3B_{01}+C_{10})/4$, so that $4r=(4,3,5,5)$; again $r\in P'$ is not a
lattice point. The candidate faces containing $r$ in their relative
interior are the $17$ weighted sets with sum $4$ and weighted sum
$(4,3,5,5)$ (Lemma~\ref{lem:enum}). Table~\ref{tab:step3} rules out each of
them by Lemma~\ref{lem:common}, using an edge from Step~1 or the tetrahedron
$F=A_{11}I_1I_2I_3$ of Step~2. Nine of them contain a segment through the
center that crosses the edge $I_2I_3$ of $F$: for example
$B_{01}+C_{10}=I_2+I_3$ rules out $3B_{01}+C_{10}$ itself. So no face of $T$
contains $r$ in its relative interior. This contradicts
Lemma~\ref{lem:den}, and proves Theorem~\ref{thm:main}. \qed

\begin{lemma}\label{lem:enum}
The sets of distinct lattice points of $P'$ with positive integer weights
summing to $4$ and weighted sum $(3,5,5,6)$ are $A_{11}+I_1+I_2+I_3$ and the
$20$ sets in the first column of Table~\ref{tab:step2}. Those with weighted
sum $(4,3,5,5)$ are the $17$ sets in the first column of
Table~\ref{tab:step3}.
\end{lemma}

The proof, a short case analysis by the last coordinates, is in
Appendix~\ref{app:enum}.

\begin{remark}
The proof uses that two faces of $T$ meet in a common face. A unimodular
cover need not satisfy this, and $P$ has one (Theorem~\ref{thm:cover}).
\end{remark}

\begin{proof}[Proof of Corollary~\ref{cor:all}]
The polytope $Q=P\times[0,1]^{d-4}$ is the lattice parallelepiped spanned
by the vectors $e_1,e_2,e_3,(1,2,3,5)$, padded with zeros, and by
$e_5,\dots,e_d$. Suppose $T$ is a unimodular
triangulation of $Q$. Its faces contained in the facet $P\times\{0\}$ of $Q$
form a triangulation of that facet (the faces of $T$ in a face of $Q$
triangulate it). Each of them is a face of a unimodular simplex, hence a
unimodular simplex for the lattice $\Z^4\times\{0\}$: its edge vectors at
a vertex are part of a basis of $\Z^d$ and lie in $\Z^4\times\{0\}$, so
they form a basis of it. This is a unimodular triangulation of $P$, contradicting
Theorem~\ref{thm:main}. Parallelepipeds are zonotopes, Minkowski sums of
segments.
\end{proof}

\section{The computer-assisted results}\label{sec:computed}

\paragraph{A second proof of Theorem~\ref{thm:main}.}
A \emph{deletion order} for a lattice polytope $Q$ lists all unimodular
simplices in $Q$ so that each one, $\sigma$, has a facet $F$ not contained
in a facet of $Q$ such that every unimodular simplex in $Q$ with facet $F$
on the other side of the hyperplane spanned by $F$ comes before $\sigma$. A
polytope with a deletion order has no unimodular triangulation
(Lemma~\ref{lem:delete}). $P$ has a deletion order of its $3380$ unimodular
simplices, which proves Theorem~\ref{thm:main} again.

\paragraph{Theorem~\ref{thm:cover}.}
A subtraction tree (Appendix~\ref{app:certs}) covers $P$ by $2441$ of its
$3380$ unimodular simplices; each of its $33{,}477$ branches ends in an
empty polyhedron.

\paragraph{Theorem~\ref{thm:classes}.}
Every lattice parallelepiped is a translate of some $P(B)$. Multiplying $B$
on the left by a matrix in $GL_4(\Z)$ brings it to Hermite normal form $H$:
upper triangular, with positive diagonal, $\det H=n$ and
$0\le H_{ij}<H_{jj}$ for $i<j$. This does not change the class of $P(B)$,
and neither does permuting the columns of $B$ and changing their signs,
which translates $P(B)$. Conversely, an equivalence maps vertices to
vertices and edges to edges, and the edge vectors at a vertex of $P(B)$ are
the columns of $B$ up to sign; so if $UP(B')+u=P(B)$, then $UB'=BD\Pi$ for
a diagonal sign matrix $D$ and a permutation matrix $\Pi$, and $B'$ and
$BD\Pi$ have the same Hermite normal form. So the classes are the Hermite
normal forms up to this relation; the table counts forms and classes. For
each $n\le 8$ a list gives one matrix per class and, for each class, either
a unimodular triangulation or a deletion order (Appendix~\ref{app:certs}).
Determinant~$1$ gives the unit cube.

\begin{center}
\begin{tabular}{lccccccc}
\toprule
determinant $n$ & 2 & 3 & 4 & 5 & 6 & 7 & 8\\
\midrule
Hermite normal forms & 15 & 40 & 155 & 156 & 600 & 400 & 1395\\
classes & 4 & 4 & 16 & 8 & 26 & 12 & 61\\
without a unimodular triangulation & 0 & 0 & 0 & 1 & 0 & 1 & 2\\
\bottomrule
\end{tabular}
\end{center}

\paragraph{Proposition~\ref{prop:cones}.}
At each vertex $v$ of $P$, a triangulation of the vertex cone into
unimodular cones with rays through lattice points of $P-v$ was found with
the SAT solver CaDiCaL~\cite{CaDiCaL} and stored;
Appendix~\ref{app:certs} describes its two checks. At $v=0$ the
Hilbert basis of the cone consists of $e_1,e_2,e_3,(1,2,3,5)$ and the four
interior lattice points of $P$; the $11$ other nonzero lattice points of
$P$, the sums of at least two of the generators, are reducible.

\paragraph{Proposition~\ref{prop:family}.}
For each odd $p$ from $5$ to $31$ a deletion order removes all unimodular
simplices of $P[1,2,2,p]$. The cases $p=5$ and $p=7$ are the classes of
determinant $5$ and $7$ in Theorem~\ref{thm:classes}, and $P[1,2,2,p]$ has a
unimodular triangulation for $1\le p\le4$ and $p=6,8$. This suggests a
question.

\begin{question}
Is it true that $P[1,2,2,p]$ has no unimodular triangulation for every odd
$p\ge5$?
\end{question}

\section{Computational verification}\label{sec:verification}

Theorem~\ref{thm:main} and Corollary~\ref{cor:all} have hand proofs. A short
program re-checks every finite statement in the proof of
Theorem~\ref{thm:main} by brute force, reading Tables~1--3 from this
paper's source.

The other results rest on certificates checked in exact integer and
rational arithmetic by programs that share no code with the searches that
produced the certificates:
\begin{itemize}
\item the deletion orders for $P$, for the four exceptional classes of
Theorem~\ref{thm:classes} and for Proposition~\ref{prop:family}, each
checked by two separately written programs, one in Python and one in C++,
that recompute the lattice points and unimodular simplices of the
parallelepiped from its generators;
\item the subtraction tree of Theorem~\ref{thm:cover}, replayed by one
program that recomputes every Farkas certificate and the unimodular
simplices of $P$ from its generators;
\item the class lists of Theorem~\ref{thm:classes}, each checked by two
separately written programs that re-enumerate all Hermite normal forms and
verify every listed triangulation, one deciding proper intersection by
circuits and the other by separating hyperplanes, proposed by linear
programming and verified exactly;
\item the vertex fans of Proposition~\ref{prop:cones}, checked by two
separately written programs that recompute the lattice points of $P$ from
its generators, one deciding proper intersection by circuits and covering
by volumes, the other by Fourier--Motzkin elimination and by neighboring
cones; both recompute the Hilbert basis at $0$.
\end{itemize}
The checkers of the hand proof, the deletion orders, the subtraction tree
and the vertex fans are also run on deliberately corrupted inputs (a wrong
identity, a missing row, a deletion order with its last step moved first, a
truncated tree, a fan with one cone dropped or replaced) and reject them.
The checks rely on the correctness of the Python and C++ implementations,
their compilers and runtimes; there is no formal proof in a proof
assistant. All checks run in about half an hour on an ordinary computer,
with under 1~GB of memory; the instructions are in the accompanying
repository.


\begin{appendices}
\section{Tables and the enumeration}\label{app:enum}

The notation is that of Step~2 in Section~\ref{sec:proof}.

\begin{table}[ht]
\centering\footnotesize
\caption{Step 1: the edge $uw$, the sum $2y=u+w$, and every split of its last coordinate.}\label{tab:step1}
% table: step 1
\begin{tabular}{lll}
\toprule
edge & $2y$ & splits of the last coordinate\\
\midrule
%% step1
$A_{00}\,I_1$ & $(1,1,1,1)$ & $0+1$: $I_1 + (0,0,0,0) = I_1 + A_{00}$ \\
$A_{01}\,I_1$ & $(1,1,2,1)$ & $0+1$: $I_1 + (0,0,1,0) = I_1 + A_{01}$ \\
$A_{10}\,I_1$ & $(1,2,1,1)$ & $0+1$: $I_1 + (0,1,0,0) = I_1 + A_{10}$ \\
$A_{11}\,I_2$ & $(1,2,2,2)$ & $0+2$: $I_2 + (0,1,1,0) = I_2 + A_{11}$; $1+1$ needs $I_1$ twice \\
$A_{00}\,I_3$ & $(1,2,2,3)$ & $0+3$: $I_3 + (0,0,0,0) = I_3 + A_{00}$; $1+2$: $I_1 + I_2 = (2,2,2,3) \neq  2y$ \\
$A_{01}\,I_3$ & $(1,2,3,3)$ & $0+3$: $I_3 + (0,0,1,0) = I_3 + A_{01}$; $1+2$: $(2,2,2,3) \neq  2y$ \\
$A_{10}\,I_3$ & $(1,3,2,3)$ & $0+3$: $I_3 + (0,1,0,0) = I_3 + A_{10}$; $1+2$: $(2,2,2,3) \neq  2y$ \\
$B_{01}\,I_3$ & $(2,2,3,3)$ & $0+3$: $I_3 + (1,0,1,0) = I_3 + B_{01}$; $1+2$: $(2,2,2,3) \neq  2y$ \\
$B_{10}\,I_3$ & $(2,3,2,3)$ & $0+3$: $I_3 + (1,1,0,0) = I_3 + B_{10}$; $1+2$: $(2,2,2,3) \neq  2y$ \\
$I_1\,C_{01}$ & $(2,3,4,6)$ & $1+5$: $I_1 + (1,2,3,5) = I_1 + C_{01}$; $2+4$: $I_2 + I_4 = (2,3,3,6) \neq  2y$; $3+3$ needs $I_3$ twice \\
$I_1\,D_{01}$ & $(3,3,4,6)$ & $1+5$: $I_1 + (2,2,3,5) = I_1 + D_{01}$; $2+4$: $(2,3,3,6) \neq  2y$; $3+3$ needs $I_3$ twice \\
$I_2\,D_{11}$ & $(3,4,4,7)$ & $2+5$: $I_2 + (2,3,3,5) = I_2 + D_{11}$; $3+4$: $I_3 + I_4 = (2,4,4,7) \neq  2y$ \\
\bottomrule
\end{tabular}
% end table
\end{table}

\begin{table}[ht]
\centering\small
\caption{Step 2: the weighted sets with sum $(3,5,5,6)$ other than $A_{11}+I_1+I_2+I_3$, a known face $R$, and an identity of Lemma~\ref{lem:common}.}\label{tab:step2}
% table: step 2
\begin{tabular}{lll}
\toprule
weighted set $S$ & face $R$ & identity\\
\midrule
$A_{00}+B_{11}+2I_3$ & $A_{11}\,I_2$ & $2A_{00}+B_{11}+2I_3=2A_{11}+3I_2$ \\
$A_{01}+B_{10}+2I_3$ & $A_{10}\,I_1$ & $A_{01}+2B_{10}+I_3=A_{10}+3I_1$ \\
$A_{10}+B_{01}+2I_3$ & $A_{01}\,I_1$ & $A_{10}+2B_{01}+I_3=A_{01}+3I_1$ \\
$A_{11}+B_{00}+2I_3$ & $A_{00}\,I_1$ & $A_{11}+2B_{00}+I_3=A_{00}+3I_1$ \\
$A_{11}+2I_1+I_4$ & $A_{00}\,I_3$ & $A_{11}+I_1+2I_4=A_{00}+3I_3$ \\
$2A_{11}+I_1+D_{00}$ & $I_2\,D_{11}$ & $A_{11}+2D_{00}+I_1=D_{11}+3I_2$ \\
$A_{00}+A_{11}+I_1+D_{11}$ & $A_{11}\,I_2$ & $2A_{00}+D_{11}+I_1=A_{11}+3I_2$ \\
$A_{00}+B_{11}+I_1+C_{11}$ & $I_3$ & $C_{11}+I_1=2I_3$ \\
$A_{01}+A_{10}+I_1+D_{11}$ & $B_{01}\,I_3$ & $A_{01}+D_{11}+I_1=B_{01}+2I_3$ \\
$A_{01}+A_{11}+I_1+D_{10}$ & $B_{10}\,I_3$ & $A_{11}+D_{10}+I_1=B_{10}+2I_3$ \\
$A_{01}+B_{10}+I_1+C_{11}$ & $I_3$ & $C_{11}+I_1=2I_3$ \\
$A_{01}+B_{11}+I_1+C_{10}$ & $A_{00}\,I_3$ & $A_{01}+C_{10}+I_1=A_{00}+2I_3$ \\
$A_{10}+A_{11}+I_1+D_{01}$ & $B_{01}\,I_3$ & $A_{11}+D_{01}+I_1=B_{01}+2I_3$ \\
$A_{10}+B_{01}+I_1+C_{11}$ & $I_3$ & $C_{11}+I_1=2I_3$ \\
$A_{10}+B_{11}+I_1+C_{01}$ & $A_{00}\,I_3$ & $A_{10}+C_{01}+I_1=A_{00}+2I_3$ \\
$A_{11}+B_{00}+I_1+C_{11}$ & $I_3$ & $C_{11}+I_1=2I_3$ \\
$A_{11}+B_{01}+I_1+C_{10}$ & $A_{10}\,I_3$ & $A_{11}+C_{10}+I_1=A_{10}+2I_3$ \\
$A_{11}+B_{10}+I_1+C_{01}$ & $A_{01}\,I_3$ & $A_{11}+C_{01}+I_1=A_{01}+2I_3$ \\
$A_{11}+B_{11}+I_1+C_{00}$ & $A_{00}\,I_3$ & $A_{11}+C_{00}+I_1=A_{00}+2I_3$ \\
$A_{11}+B_{11}+I_2+I_4$ & $I_1$ & $B_{11}+I_2=2I_1$ \\
\bottomrule
\end{tabular}
% end table
\end{table}

\begin{table}[ht]
\centering\small
\caption{Step 3: the weighted sets with sum $(4,3,5,5)$, a known face $R$, and an identity of Lemma~\ref{lem:common}.}\label{tab:step3}
% table: step 3
\begin{tabular}{lll}
\toprule
weighted set $S$ & face $R$ & identity\\
\midrule
$3B_{01}+C_{10}$ & $A_{11}\,I_1\,I_2\,I_3$ & $B_{01}+C_{10}=I_2+I_3$ \\
$A_{00}+2B_{01}+D_{11}$ & $A_{11}\,I_1\,I_2\,I_3$ & $A_{00}+D_{11}=I_2+I_3$ \\
$A_{01}+2B_{01}+D_{10}$ & $A_{11}\,I_1\,I_2\,I_3$ & $A_{01}+D_{10}=I_2+I_3$ \\
$A_{10}+2B_{01}+D_{01}$ & $A_{11}\,I_1\,I_2\,I_3$ & $A_{10}+D_{01}=I_2+I_3$ \\
$A_{11}+2B_{01}+D_{00}$ & $A_{11}\,I_1\,I_2\,I_3$ & $A_{11}+D_{00}=I_2+I_3$ \\
$B_{00}+2B_{01}+C_{11}$ & $A_{11}\,I_1\,I_2\,I_3$ & $B_{00}+C_{11}=I_2+I_3$ \\
$2B_{01}+B_{10}+C_{01}$ & $A_{11}\,I_1\,I_2\,I_3$ & $B_{10}+C_{01}=I_2+I_3$ \\
$2B_{01}+B_{11}+C_{00}$ & $A_{11}\,I_1\,I_2\,I_3$ & $B_{11}+C_{00}=I_2+I_3$ \\
$2B_{01}+I_1+I_4$ & $A_{11}\,I_1\,I_2\,I_3$ & $I_1+I_4=I_2+I_3$ \\
$2B_{01}+I_2+I_3$ & $I_1\,C_{01}$ & $B_{01}+I_2+2I_3=C_{01}+3I_1$ \\
$A_{00}+B_{01}+B_{11}+D_{01}$ & $A_{01}\,I_1$ & $2A_{00}+3B_{11}+D_{01}=A_{01}+5I_1$ \\
$A_{01}+B_{00}+B_{01}+D_{11}$ & $I_1\,D_{01}$ & $A_{01}+2B_{00}+B_{01}+3D_{11}=2D_{01}+5I_1$ \\
$A_{01}+B_{00}+B_{11}+D_{01}$ & $B_{01}\,I_3$ & $3A_{01}+B_{00}+4B_{11}+3D_{01}=6B_{01}+5I_3$ \\
$A_{01}+B_{01}+B_{10}+D_{01}$ & $A_{11}\,I_1\,I_2\,I_3$ & $A_{01}+2B_{10}+D_{01}=3I_1+I_2$ \\
$A_{01}+B_{01}+B_{11}+D_{00}$ & $I_1\,D_{01}$ & $A_{01}+3B_{11}+2D_{00}=D_{01}+5I_1$ \\
$A_{11}+B_{00}+B_{01}+D_{01}$ & $A_{01}\,I_1$ & $3A_{11}+2B_{00}+B_{01}+D_{01}=2A_{01}+5I_1$ \\
$B_{00}+B_{01}+B_{11}+C_{01}$ & $B_{01}\,I_3$ & $B_{00}+4B_{11}+3C_{01}=3B_{01}+5I_3$ \\
\bottomrule
\end{tabular}
% end table
\end{table}

\begin{proof}[Proof of Lemma~\ref{lem:enum}]
The lattice points with last coordinate $0$ are the $A_{ab}$ and $B_{ab}$,
with last coordinate $5$ the $C_{ab}$ and $D_{ab}$; the first coordinate is
$0$ on $A_{ab}$, $1$ on $B_{ab}$, $C_{ab}$ and every $I_k$, and $2$ on
$D_{ab}$.

\emph{Sum $(3,5,5,6)$.} Let $w$ be the total weight on $C$ and $D$ points.
The last coordinates sum to $6$, so $w\le1$.
If $w=1$, the interior points contribute $1$ to the last coordinate, so
$I_1$ has weight $1$ and weight $2$ lies on $A$ and $B$ points. With a
point $D_{ab}$ the weight-$2$ part has first coordinate $0$, so it consists
of $A$ points with $x_2$-sum $2-a$ and $x_3$-sum $2-b$: this gives
$2A_{11}$ with $D_{00}$; $\{A_{10},A_{11}\}$ with $D_{01}$; $\{A_{01},A_{11}\}$
with $D_{10}$; and $\{A_{00},A_{11}\}$ or $\{A_{01},A_{10}\}$ with $D_{11}$,
five sets. With a point $C_{ab}$ the weight-$2$ part has first coordinate
$1$, so it is one $A$ and one $B$ point, with $x_2$-sum $2-a$ and $x_3$-sum
$2-b$; this gives $1,2,2,4$ sets for $C_{00},C_{01},C_{10},C_{11}$, nine sets.
If $w=0$, the interior weights contribute $6$ to the last coordinate. The
first coordinates sum to $3$, which is the interior weight plus the
$B$ weight, so exactly one $A$ point occurs, with weight $1$. The interior
parts that leave room for it are $I_4+I_2$, $I_4+2I_1$, $2I_3$, $I_3+I_2+I_1$
and $3I_2$. Matching the second and third coordinates (both $5$) gives
$A_{11}+B_{11}+I_2+I_4$; $A_{11}+2I_1+I_4$; the four sets
$A_{ab}+B_{a'b'}+2I_3$ with $a+a'=b+b'=1$; $A_{11}+I_1+I_2+I_3$; and nothing
for $3I_2$, which would need $a=2$. That is seven sets, and $21$ in all.

\emph{Sum $(4,3,5,5)$.} The last coordinates sum to $5$, so the total weight
on $C$ and $D$ points is at most $1$. With one point $C_{ab}$, the weight-$3$
rest has first coordinate $3$, so it consists of $B$ points, with $x_2$-sum
$1-a$ and $x_3$-sum $3-b$: $C_{00}$ gives $2B_{01}+B_{11}$; $C_{10}$ gives
$3B_{01}$; $C_{01}$ gives $2B_{01}+B_{10}$ and $B_{00}+B_{01}+B_{11}$; $C_{11}$
gives $B_{00}+2B_{01}$; five sets. With one point $D_{ab}$, the rest has
first coordinate $2$, so it is one $A$ point and $B$ weight $2$, with
$x_2$-sum $1-a$ and $x_3$-sum $3-b$: $D_{00}$ gives two sets, $D_{10}$ one,
$D_{01}$ five and $D_{11}$ two, ten sets in all, as listed in
Table~\ref{tab:step3}. With no $C$ or $D$ point, the interior weights
contribute $5$ to the last coordinate and the first coordinates force the
interior weight plus the $B$ weight to be $4$, so no $A$ point occurs. The
interior parts are $I_4+I_1$, $I_3+I_2$, $I_3+2I_1$, $2I_2+I_1$ and
$I_2+3I_1$; the second coordinate $3$ excludes $I_3+2I_1$ and $I_2+3I_1$,
$2I_2+I_1$ would need a $B$ point with third coordinate $2$, and $I_4+I_1$
and $I_3+I_2$ each need $2B_{01}$. That is two sets, and $17$ in all.
\end{proof}

\clearpage
\section{The certificates}\label{app:certs}

\paragraph{Deletion orders.}
\begin{lemma}\label{lem:delete}
Let $T$ be a unimodular triangulation of a lattice polytope $Q$, let
$\sigma\in T$, and let $F$ be a facet of $\sigma$ not contained in a facet
of $Q$. Then some $\tau\in T$ has $F$ as a facet and lies on the other side
of the hyperplane spanned by $F$. Hence a polytope with a deletion order has
no unimodular triangulation.
\end{lemma}

\begin{proof}
A point $x\in\relint F$ lies in the interior of $Q$, since otherwise a
supporting hyperplane of $Q$ through $x$ would contain $F$, and $F$ would
lie in a facet of $Q$. Points near $x$ on the other side of $F$ lie in $Q$
and are covered by the finitely many simplices of $T$, so one of them,
$\tau$, contains such points arbitrarily close to $x$, and $x$ itself. Then
$\sigma\cap\tau$ is a common face containing $x$, hence $F$; it is not
$\sigma$, since $\tau$ has points on the other side, so $F$ is a facet of
$\tau$, and $\tau$ lies on the other side. If $Q$ had a deletion order and
a unimodular triangulation $T\ne\emptyset$, the first simplex of $T$ in the
order would have such a $\tau$ before it, a contradiction.
\end{proof}

\paragraph{The subtraction tree.}
We cover the interior of $P$ recursively. Given an open polyhedron $R$ and a
unimodular simplex $S$ in $P$ with facet inequalities $f_0,\dots,f_4\ge0$,
the part of $R$ outside $S$ is, up to finitely many hyperplanes, the union of
the open polyhedra $R\cap\{f_0>0,\dots,f_{i-1}>0,\ f_i<0\}$ for
$i=0,\dots,4$; we recurse on each. A branch stops when its polyhedron
$\{x: a_k\cdot x<b_k \text{ for all } k\}$ is empty, which an exact Farkas
certificate proves: numbers $y_k\ge0$, not all zero, with
$\sum_k y_ka_k=0$ and $\sum_k y_kb_k\le0$. The tree for $P$ starts from its
interior, so the interior of $P$ minus the union of the subtracted simplices
lies in finitely many hyperplanes; the union, which is closed, is $P$.

\paragraph{Triangulations of the classes.}
A list of $24n=4!\cdot n$ distinct unimodular simplices, the normalized
volume of $P(B)$, with vertices among the lattice points of $P(B)$, is a
triangulation when any two meet in a common face: their interiors are then
disjoint, their volumes add up to that of $P(B)$, and the closed union is
$P(B)$. One checker decides this by circuits: two simplices fail to meet in
a common face exactly when some minimal affinely dependent set of their
vertices has the points with positive coefficients in one simplex and those
with negative coefficients in the other. The other checker verifies, in
exact arithmetic, a hyperplane through their common vertices that separates
them, found by linear programming.

\paragraph{Vertex fans.}
Let $v$ be a vertex of $P$, let $g_1,\dots,g_4$ be the edge vectors of $P$
at $v$ (the columns of $B$ up to sign), so that the vertex cone is
$C_v=\operatorname{cone}(g_1,\dots,g_4)$, and let $h(x)$ be the sum of the
coordinates of $x$ in the basis $g_1,\dots,g_4$. Each stored cone is spanned
by $p_1-v,\dots,p_4-v$ for lattice points $p_i$ of $P$ with
$|\det(p_1-v,\dots,p_4-v)|=1$, and lies in $C_v$. Both checkers verify this
from the lattice points of $P$, recomputed from its generators. The first
checks that the cones meet pairwise in common faces, by the circuit test on
the simplices $\conv(0,p_1-v,\dots,p_4-v)$, and that
\[
\sum\frac{1}{h(p_1-v)\,h(p_2-v)\,h(p_3-v)\,h(p_4-v)}=|\det B|,
\]
summed over the cones. Up to one common factor, the left side is the total
volume of the sections of the cones by the hyperplane $h=1$, and the right
side is the volume of the section of $C_v$. With disjoint interiors, the
closed sections then fill the section of $C_v$, so the cones fill $C_v$.
The second checker shows by Fourier--Motzkin elimination that no point of
two cones has a positive coordinate, in the basis of one of them, at a ray
the other lacks; so the cones meet in common faces. It then checks that
every facet of a cone that is not contained in a facet of $C_v$ is a facet of
another cone on its other side. Then the cones fill $C_v$: otherwise a
generic segment from the interior of a cone to an interior point of $C_v$
outside their union would leave the union through such a facet.
Both checkers compute the Hilbert basis of $C_0$ from the lattice points of
$P$: its elements lie in $P$, and a lattice point of $P$ is reducible exactly
when it is the sum of two nonzero lattice points of $P$.

\section{Sources not read in full}\label{app:unread}

The thesis of von Thaden~\cite{vT}, the paper~\cite{BG} and the published
edition of the book~\cite{BGbook} of Bruns and Gubeladze (we read a 2006
draft of the book) are all cited in~\cite{HPPS}, whose authors posed the
question after them. Sebő's paper~\cite{Se} is cited only for its theorem
on 3-dimensional cones, as reported in~\cite{BGT}. For Blanco's
chapter~\cite{Bl} we read its title and its zbMATH summary (Zbl
1428.52019), which is about covers: ``In this paper I will study a
3-parameter family of lattice parallelepipeds. I will study the behavior of
the lattice with respect to the parameters and prove that three 2-parameter
subfamilies have unimodular covers.''
\end{appendices}

\begin{thebibliography}{99}
\bibitem{CaDiCaL} A.~Biere, T.~Faller, K.~Fazekas, M.~Fleury, N.~Froleyks
and F.~Pollitt, CaDiCaL 2.0, in \emph{Computer Aided Verification (CAV
2024)}, Lecture Notes in Comput.\ Sci.\ 14681, Springer, 2024, 133--152,
\href{https://doi.org/10.1007/978-3-031-65627-9_7}{doi:10.1007/978-3-031-65627-9\_7}.

\bibitem{Bl} M.~Blanco, Some lattice parallelepipeds with unimodular covers,
in \emph{Algebraic and Geometric Combinatorics on Lattice Polytopes}, World
Scientific, 2019, 85--100,
\href{https://doi.org/10.1142/9789811200489_0005}{doi:10.1142/9789811200489\_0005}.

\bibitem{BGS} C.~Bouvier and G.~Gonzalez-Sprinberg, Système générateur
minimal, diviseurs essentiels et $G$-désingularisations de variétés
toriques, \emph{Tôhoku Math.\ J.} 47 (1995), no.~1, 125--149,
\href{https://doi.org/10.2748/tmj/1178225640}{doi:10.2748/tmj/1178225640}.

\bibitem{BG} W.~Bruns and J.~Gubeladze, Normality and covering properties of
affine semigroups, \emph{J. Reine Angew.\ Math.} 510 (1999), 161--178,
\href{https://doi.org/10.1515/crll.1999.044}{doi:10.1515/crll.1999.044}.

\bibitem{BGbook} W.~Bruns and J.~Gubeladze, \emph{Polytopes, Rings, and
K-Theory}, Springer Monographs in Mathematics, Springer, 2009,
\href{https://doi.org/10.1007/b105283}{doi:10.1007/b105283}.

\bibitem{BGT97} W.~Bruns, J.~Gubeladze and N.~V.~Trung, Normal polytopes,
triangulations, and Koszul algebras, \emph{J. Reine Angew.\ Math.} 485
(1997), 123--160,
\href{https://doi.org/10.1515/crll.1997.485.123}{doi:10.1515/crll.1997.485.123}.

\bibitem{BGT} W.~Bruns, J.~Gubeladze and N.~V.~Trung, Problems and
algorithms for affine semigroups, \emph{Semigroup Forum} 64 (2002),
180--212,
\href{https://doi.org/10.1007/s002330010099}{doi:10.1007/s002330010099};
quoted from \href{https://arxiv.org/abs/math/0101116}{arXiv:math/0101116}.

\bibitem{CS} G.~Codenotti and F.~Santos, Unimodular covers of
$3$-dimensional parallelepipeds and Cayley sums, \emph{Combinatorial
Theory} 3 (2023), no.~3,
\href{https://doi.org/10.5070/C63362785}{doi:10.5070/C63362785}; numbering
and quotations from
\href{https://arxiv.org/abs/1907.12312v1}{arXiv:1907.12312v1}.

\bibitem{FH} L.~Ferroni and A.~Higashitani, Examples and counterexamples in
Ehrhart theory, \emph{EMS Surv.\ Math.\ Sci.} (2024),
\href{https://doi.org/10.4171/emss/86}{doi:10.4171/emss/86}; section
numbers from \href{https://arxiv.org/abs/2307.10852v4}{arXiv:2307.10852v4}.

\bibitem{FZ} R.~T.~Firla and G.~M.~Ziegler, Hilbert bases, unimodular
triangulations, and binary covers of rational polyhedral cones,
\emph{Discrete Comput.\ Geom.} 21 (1999), 205--216,
\href{https://doi.org/10.1007/PL00009416}{doi:10.1007/PL00009416}; quoted
from the authors' revised preprint of September 24, 1997.

\bibitem{HPPS} C.~Haase, A.~Paffenholz, L.~C.~Piechnik and F.~Santos,
Existence of unimodular triangulations---positive results, \emph{Mem.\
Amer.\ Math.\ Soc.} 270 (2021), no.~1321,
\href{https://doi.org/10.1090/memo/1321}{doi:10.1090/memo/1321}; page and
section numbers from
\href{https://arxiv.org/abs/1405.1687v3}{arXiv:1405.1687v3}.

\bibitem{Se} A.~Sebő, Hilbert bases, Carathéodory's theorem and
combinatorial optimization, in \emph{Integer Programming and Combinatorial
Optimization} (R.~Kannan and W.~R.~Pulleyblank, eds.), University of
Waterloo Press, 1990, 431--456.

\bibitem{vT} M.~von Thaden, Unimodular covers and triangulations of lattice
polytopes, Ph.D. thesis, Universität Osnabrück, 2007.
\end{thebibliography}
\end{document}
